% Provenance: copied from the design workspace CambridgePQ/recursion.fast, `report/sections/09-recursion.tex`, on 2026-10-19. Read-only copy; the source is authoritative.
% 09-recursion -- written by W3
\section{Recursion and the batcher}\label{sec:recursion}

The brief asks for a two-to-one \emph{batcher}: a program that turns two proofs into one, so that a
tree of batch steps compresses many proofs into a single one. This section specifies the batcher as
an accumulation scheme in the information-theoretic (IT) model, the model of
\cref{sec:protocol,sec:security}. It then says what composing batch steps proves, what cannot be a
Lean theorem at all, and which variants are deferred. Background on IVC, PCD, accumulation and the
negative results is in \cref{subsec:bg-recursion}; bracketed tags such as [A2-M3] are explained in
\cref{sec:summary}. Terms used throughout:
\begin{itemize}[nosep]
\item The \emph{inner verifier} is the verifier that recursion runs inside the next proof.
\item The \emph{final verifier} (or \emph{decider}) runs once, at the root, usually natively.
\item An \emph{accumulator} is a statement--witness pair of a relation $\Racc$. Its statement is the
  only state that one batch step hands to the next.
\end{itemize}

%-----------------------------------------------------------------------------------------------
\subsection{What is wrong with the brief's definition}\label{subsec:rec-critique}

The brief defines the batcher by one property: \enquote{if $\Vrf$ accepts $\pi_1$ and $\pi_2$, then
$\Vrf$ accepts $\pi'$ with high probability}. That is a \emph{completeness} property, and two useless
programs satisfy it.
\begin{itemize}[nosep]
\item \emph{Constant proof.} If some fixed $\pi_0$ is accepted, the batcher $(\pi_1,\pi_2)\mapsto\pi_0$
  meets the definition. Acceptance of $\pi'$ then says nothing about $\pi_1$ or $\pi_2$.
\item \emph{Concatenation.} With $\pi'=(\pi_1,\pi_2)$ and a final verifier that checks both halves,
  the batcher is sound and complete. It also compresses nothing: after a tree of depth $d$ the proof
  has $2^d$ parts.
\end{itemize}
The definition leaves three further questions open.
\begin{enumerate}[nosep,label=(\roman*)]
\item \emph{Soundness of the batched object.} The property that carries security is \emph{knowledge
  soundness}: from any prover that makes $\pi'$ accept one must extract witnesses for \emph{both}
  input statements.
\item \emph{What the batcher sees.} If it sees finished proofs only, it is \emph{aggregation}
  (SnarkPack-style). If it also sees instance--witness pairs or prover state, it is
  \emph{accumulation} or \emph{folding}. Only accumulation keeps the recursion circuit small,
  because that circuit runs the cheap accumulation verifier and never the decider
  \cite{BCLMS21,KST22-Nova}.
\item \emph{Inner versus final verifier.} The brief allows the inner and final verifiers to differ,
  but it does not say which one is scored. Recursion pays the inner one at every node; the final
  one runs once.
\end{enumerate}

\begin{decision}{REC-1}{The batcher is a split accumulation scheme in the IT model}
A T2 submission consists of three interactive reductions of knowledge
(\cref{def:bg-rok}): $\mathsf{embed}:\Rrcs\to\Racc$, $\mathsf{fold}:\Racc\times\Racc\to\Racc$ and
$\mathsf{decide}:\Racc\to\top$. Each reduction must satisfy the base track's certificate. That
certificate has four parts: perfect completeness for every challenge vector; worst-case round-by-round
knowledge soundness (RBR-KS) with a named extractor, in counting form; the per-round error cap
$\eps_i\le2^{-128}$; and the round cap $\le2^{14}$. The reductions use the same contract-fixed
round shapes (\cref{sec:protocol}). The accumulator statement has a \emph{contract-fixed shape per
class} (\cref{subsec:rec-shape}). \emph{Rationale:} this is the reduction-of-knowledge form that
compiles, via BCS and Fiat--Shamir, into a non-interactive split accumulation scheme with
straight-line extraction \cite[Thms.~B.4 and 10.3]{BCFW25-WARP}. It is also leanVM's deployed design:
\enquote{deferred claims} on fixed polynomials, discharged at the root \cite{leanVM}. \emph{Answers:}
the constant-proof and concatenation batchers; an unscored or ambiguous inner verifier.
[DESIGN §6; D1 §4.2]
\end{decision}

%-----------------------------------------------------------------------------------------------
\subsection{Definition}\label{subsec:rec-def}

Fix a class $k\in\{16,20,24\}$ with its class shape (\cref{sec:protocol}). Let $\mathsf{St}_k$ be
the base track's statement type: the verification key $\vk_i$ and the public input $x$, with
$\idx(i)$ as oracle~0. Let $\mathsf{AccSt}_k$ be the accumulator statement type fixed by the shape
of \cref{subsec:rec-shape}, and let $\mathsf{AccW}$ be a submitter-chosen witness type.

\begin{definition}[Accumulation scheme, IT model]\label{def:rec-acc}
An \emph{accumulation scheme} for class $k$ consists of a relation
$\Racc\subseteq\mathsf{AccSt}_k\times\mathsf{AccW}$ and three public-coin reductions:
\begin{align*}
  \mathsf{embed} &: (\mathsf{St}_k,\ \mathsf{Wit}_k)\to(\mathsf{AccSt}_k,\ \mathsf{AccW})
     &&\text{from } \Rrcs^{(k)} \text{ to } \Racc,\\
  \mathsf{fold} &: (\mathsf{AccSt}_k^2,\ \mathsf{AccW}^2)\to(\mathsf{AccSt}_k,\ \mathsf{AccW})
     &&\text{from } \mathsf{pair}(\Racc) \text{ to } \Racc,\\
  \mathsf{decide} &: (\mathsf{AccSt}_k,\ \mathsf{AccW})\to(\mathsf{Unit},\mathsf{Unit})
     &&\text{from } \Racc \text{ to } \top,
\end{align*}
where $\mathsf{pair}(\R)=\{((s_1,s_2),(w_1,w_2)) : (s_1,w_1)\in\R\wedge(s_2,w_2)\in\R\}$ and
$\top=\mathsf{Unit}\times\mathsf{Unit}$. The scheme is \emph{secure} if each reduction is perfectly
complete and RBR knowledge sound for its pair of relations, with every per-round error at most
$2^{-128}$ and at most $2^{14}$ rounds.
\end{definition}

$\Racc$ is the submitter's choice, and that freedom is harmless. The leaves are fixed to the
contract's $\Rrcs$ and the root to $\top$. A degenerate $\Racc$ only makes $\mathsf{embed}$ or
$\mathsf{decide}$ impossible to prove. The knowledge requirement for $\mathsf{fold}$ is the one
missing from the brief: from a prover that makes the fold verifier output an accumulator in $\Racc$,
the extractor recovers $\Racc$-witnesses for \emph{both} inputs. The listing is an abridged excerpt
of the compiled \code{Accumulation.lean} (full listing in \cref{sec:appendix-contract}). A
\lean{Reduction} is a protocol whose spec verifier outputs the plain part of an accumulator and
whose output oracles are a declared selection \lean{select} among the step's QUERY oracles; its
security clause is the base track's \lean{SecureWith}.

\begin{leancode}[fontsize=\footnotesize,label={\footnotesize Contract/Accumulation.lean (abridged)}]
/-- Cap on the plain part of an accumulator, in K-elements. -/
def accCap : ℕ := 2 ^ 12
/-- Cap on the number of accumulator oracles. -/
def accOracleCap : ℕ := 16

/-- Declared accumulator shape. -/
structure AccShape where
  plain : ℕ
  oracles : List OShape

/-- An accumulator statement: plain part and one oracle of each declared shape. -/
abbrev AccStmt (A : AccShape) := (Fin A.plain → K) ×
    ((j : Fin A.oracles.length) → OVal (A.oracles.get j))

-- …

/-- **Agreement with output**: on every input and transcript the image accepts iff the spec
does, and then leaves the spec's plain output at `OUT_BASE`. -/
def Agrees (img : Image) : Prop :=
  ∀ s tr, (R.runOn view img s tr).map (fun r => readOut r.2 B.plain) =
    (R.verify s tr).map wordsOf

-- …

/-- An accumulation scheme for size `sz`. -/
structure Accumulation (sz : Sizes) where
  acc : AccShape
  acc_le : acc.plain ≤ accCap ∧ acc.oracles.length ≤ accOracleCap
  -- …
  AccWit : Type
  Racc : Set (AccStmt acc × AccWit)
  embed : Reduction (Index sz × PubInput sz) (Wit sz) AccWit acc
  fold : Reduction (AccStmt acc × AccStmt acc) (AccWit × AccWit) AccWit acc
  decide : Reduction (AccStmt acc) AccWit Unit ⟨0, []⟩

-- …

/-- **Security of the scheme**: each step secure for its relations; selections well-shaped. -/
structure Secure : Prop where
  embed : A.embed.Secure A.embedView (rel sz) A.Racc
  fold : A.fold.Secure A.foldView (pairRel A.Racc) A.Racc
  decide : A.decide.Secure A.decideView A.Racc Set.univ
  select : A.embed.SelectOk [A.idxShape] ∧ A.fold.SelectOk (A.acc.oracles ++ A.acc.oracles) ∧
    A.decide.SelectOk A.acc.oracles
\end{leancode}
\noindent The omitted fields are the holographic preprocessing of \lean{embed} (\lean{VK},
\lean{IDX}, as in \lean{Protocol}); \lean{embedView}, \lean{foldView} and \lean{decideView} fix
what each machine sees on its tape and as QUERY oracles.

%-----------------------------------------------------------------------------------------------
\subsection{The fixed accumulator shape, and why the cap must be in the definitions}\label{subsec:rec-shape}

The accumulator statement has two parts, both fixed per class.
\begin{itemize}[nosep]
\item A \emph{plain part} of at most $\mathsf{accCap}=2^{12}$ $K$-elements (32\,KiB). It is written
  on the tape and read like the public input.
\item Up to $r\le16$ \emph{accumulator oracles} of declared lengths. They are committed vectors, so
  the verifier reaches them only through the priced \texttt{QUERY} syscall.
\end{itemize}
Three interface rules complete the shape. First, $\mathsf{embed}$ uses the base track's holographic
interface: it reads $\vk_i$ and $x$ from the tape and $\idx(i)$ by \texttt{QUERY}, never the index
as a Lean value. Second, each output oracle of $\mathsf{fold}$ is a declared \emph{selection}: either
an input-accumulator oracle or a fold message. This is the discipline of ArkLib's output-oracle
embeddings, and it keeps the output oracles committed objects rather than verifier-computed ones.
Third, the machine--specification agreement $\mathsf{Agrees}$ (\cref{subsec:proto-agrees}) includes
equality of the output plain part, not only of the accept bit. A decision-only agreement would leave
the fold machine free to output a different accumulator from the one the specification proves secure.

The cap is constant; it does not grow with $n_x$. The public input enters only at $\mathsf{embed}$,
which must reduce $x$ to an evaluation claim. A cap of $2^{12}+n_x$ (D1's first proposal) would help
only designs that carry one $x$ verbatim, and no fold produces that. For instance, $x_1+r\,x_2$ with
$r\in E$ already has $3n_x$ $K$-elements [A1-m4].

\paragraph{Why the cap must be a type, not a field.} In D1's first version the accumulation structure
had a serializer $\mathsf{serAcc}:\mathsf{AccStmt}\to(\mathsf{Fin}\ \mathsf{accLen}\to F)$ and a
proof of $\mathsf{accLen}\le\mathsf{accCap}$. But $\mathsf{Accumulation.Secure}$ never mentioned
either field, and $\mathsf{serAcc}$ was unconstrained. Review A2 built the counterexample in
Lean:\footnote{\path{probes/a2-probes/A2Probes.lean}, section C,
\lean{trivialAcc0\_secure} and \lean{fullRecAcc\_secure}; sorry-free, axioms
\lean{propext}, \lean{Classical.choice}, \lean{Quot.sound}.} \lean{trivialAcc0} sets
\lean{AccStmt := Index F sz × PubInput F sz}, \lean{accLen := 0} and a vacuous \lean{serAcc}, and
\lean{trivialAcc0\_secure} proves it secure.
The accumulator \enquote{of length 0} is a whole R1CS instance, and the fold verifier reads it as a
free Lean value. A cap that the security predicate does not read is decoration. This is leanerVM's
lesson A5 again: every size cap the verifier relies on must be in the relation or an explicit
hypothesis.\footnote{\repo{leanerVM@8aaf8a0}, adversarial-review lesson A5, as summarised in
\path{notes/R5-leanervm.md} §3.3, item 4.} In v0.2 the statement \emph{type} is $\mathsf{AccStmt}\ \mathit{sh}$,
built from raw vectors whose lengths are bounded in $\mathsf{AccShape}$. A submitter can no longer
choose a type that hides data. Concatenation is impossible because the output has the same shape as
each input.

\paragraph{The oracle lengths must be typed too.} Review A4 of the compiled v0.2 contract found the
same mistake one level down [A4-F1]. That contract typed the plain part but kept the accumulator
oracles as an untyped list, \lean{AccStmt A := (Fin A.plain → K) × List Oracle}. Its fold and decide
cost clauses quantified over \emph{all} statements, including ones whose oracles do not have the
declared lengths. Take any accepting run that queries input oracle $j$, and zero-extend that oracle by
$d$ levels. Every queried leaf index is still in range, so the run still accepts with the same final
state. Only $q^{\mathrm{lev}}$ grows, by $d$ per query, so no claim bounds the accepting runs. The
probe proves this in Lean, sorry-free: a six-instruction image \lean{qimg} that queries reaches
any $q^{\mathrm{lev}}=L$ (\lean{qimg\_levels}), so \lean{decide\_cost\_unsat} and
\lean{fold\_cost\_unsat} refute every cost claim and \lean{no\_certificate2} every T2 certificate
with that decider image.
So no T2 certificate could exist for a fold or a decider that reads an accumulator oracle. Accumulators
collapse to their plain part, which in the IT model cannot bind a witness, and T2 degenerates to
\lean{fullRecAcc}. The v0.3 fix types the oracles by the declared shape, as in the listing of
\cref{subsec:rec-def}: oracle $j$ of an \lean{AccStmt A} is an \lean{OVal (A.oracles.get j)}. The
machine-visible length then equals the charged length,
and $\mathsf{Agrees}$, $\mathsf{Secure}$ and the cost clauses range over well-shaped inputs only. A
positive control, an oracle-reading decider that meets a cost bound, guards against regression.
[DESIGN §15.2; A4-F1]

\begin{decision}{REC-2}{Accumulator shape fixed per class, cap enforced by the type}
Plain part at most $2^{12}$ $K$-elements on the tape; at most $16$ accumulator oracles whose declared
lengths are part of the type, reached only by \texttt{QUERY}; holographic $\mathsf{embed}$; fold output oracles are a
declared selection; $\mathsf{Agrees}$ includes output equality. \emph{Rationale:} without a
load-bearing cap, concatenation is a sound and complete batcher, and the fold verifier gets a free
lookup table (the ots.golf lesson, \cref{subsec:prior-lessons}). \emph{Answers:} A2-M3 (Lean
counterexample \lean{trivialAcc0\_secure}); A4-F1 (Lean counterexample \lean{no\_certificate2});
A1-m4. [DESIGN §6, §15.2]
\end{decision}

%-----------------------------------------------------------------------------------------------
\subsection{Inner versus final verifier, and what is scored}\label{subsec:rec-score}

A binary tree over $N$ leaves runs $N$ embeds, $N-1$ folds and one decide. Per leaf, that is one leaf
step, about one fold step and $1/N$ of a decider. Each image is scored as on the recursion board
(\cref{subsec:cost-zk}), under each of the four zkVM profiles $P$. So $S_{P,k}$ includes the
Fiat--Shamir surcharge, the byte charge and the blow-up factor.
\[
  \mathrm{Score}_{T2}=\Bigl(\prod_{P}\prod_{k\in\{16,20,24\}}\bigl(S_{\mathrm{leaf},P,k}+S_{\mathrm{fold},P,k}
     +S_{\mathrm{dec},P,k}/64\bigr)\Bigr)^{1/12}.
\]
The constant $64$ is the amortisation of a tree with 64 leaves. Block-level aggregation uses
$2^4$--$2^{10}$ leaves, so 64 is a judgement (open item O12). There is no separate decider cap:
the decider is bounded only by the fuel $F=2^{36}$ per hart (an all-runs bound for deciders is
planned, not in \lean{Certificate2}), so linear (Arc/WARP-style) deciders are
admissible at every class and compete through the $/64$ weight [DESIGN §15.3; A5-M8]. A succinct
decider contributes 1--3\% of the score. A linear decider at $2^{24}$ costs about
$400\cdot2^{24}\approx6.7\times10^9$ cycles, so it contributes $\approx10^8$ after division and
dominates. Between $/64$ and $/1024$ the ranking of a linear-decider design against a succinct one
flips only if the succinct design's fold is at least 30\% more expensive [A1-m5]. The weight is
therefore a strong preference for succinct deciders, not a ban.
T1 is the special case with a null fold. Records are labelled \enquote{one-step} until the tree
theorem of \cref{subsec:rec-compose} exists.

\paragraph{Two baselines.} Both satisfy \cref{def:rec-acc} in Lean today [A2 §C].
\begin{itemize}[nosep]
\item \lean{fullRecAcc}: $\Racc=\mathsf{univ}$ over $\mathsf{Unit}$; $\mathsf{embed}$ is the whole
  base-track verifier; $\mathsf{fold}$ and $\mathsf{decide}$ accept everything in zero rounds. It
  meets every cap, because the accumulator is $\mathsf{Unit}$, and it scores
  $S_{\mathrm{leaf}}+0+0$, exactly the T1 score. This is not an exploit. It is what deploying full
  verification at every leaf costs, and it is the number a T2 entry must beat. A T2 entry is
  interesting only if $S_{\mathrm{leaf}}+S_{\mathrm{fold}}+S_{\mathrm{dec}}/64$ is below the best T1
  score.
\item \lean{trivialFold}: $\Racc=\Rrcs$, $\mathsf{embed}$ is the identity, and $\mathsf{fold}$ lets
  the prover send $(w_1,w_2)$, checks both instances and outputs a fixed satisfiable instance. Its
  fold cost is two full verifications. Its accumulator contains the index, so it violates the fixed
  accumulator shape. It is kept only as a non-holographic illustration of a non-vacuous certificate.
\end{itemize}

%-----------------------------------------------------------------------------------------------
\subsection{What composition gives}\label{subsec:rec-compose}

\paragraph{The tree as one IOP.} For $N=2^d$ leaves, sequentialise the tree:
$\mathsf{embed}^{N}\gg\mathsf{fold}\text{-layer}_d\gg\dots\gg\mathsf{fold}\text{-layer}_1\gg
\mathsf{decide}$. The result is a single public-coin protocol from $\Rrcs^{N}$ to $\top$. In the IT
model, sequential composition preserves perfect completeness and worst-case RBR-KS. The per-round
errors are \emph{concatenated}, not summed, and the extractor is the composition of the backward
extractors. leanerVM has proved this for guarded-first \lean{append} (582 lines,
\path{ToArkLib/KnowledgeAppend.lean}, a port of ArkLib PR~\#615); ArkLib's own
\lean{append\_rbrKnowledgeSoundness} is still \code{sorry}
\cite{arklib2026}.\footnote{\repo{leanerVM@8aaf8a0}, \path{ToArkLib/KnowledgeAppend.lean};
\repo{ArkLib@db7005f}.} A fold layer of
$N/2^j$ independent folds is a \emph{parallel} composition. It needs a frame rule (ArkLib's
\lean{liftContext} for RBR-KS, also \code{sorry}).

\begin{leancode}
-- Organiser obligation (level L2), statement only. Proof needs the frame rule.
theorem tree_secure (A : Accumulation k) (hA : A.Secure) (d : ℕ) (hd : d ≤ dMax) :
    (tree A d).Secure (relPow k (2 ^ d)) Set.univ
-- per-round errors: those of the components (cap unchanged)
-- rounds: 2^d·n_embed + (2^d − 1)·n_fold + n_decide
-- extractor: the composed backward chain
\end{leancode}

\paragraph{Error accounting.} Every per-round error stays $\le2^{-128}$. The number of rounds is
$k_{\mathrm{tree}}=N\,k_{\mathrm{embed}}+(N-1)\,k_{\mathrm{fold}}+k_{\mathrm{dec}}$, which is above
the per-reduction cap of $2^{14}$. The tree theorem therefore states its own round count. For example,
$N=2^{10}$ and 210 rounds per reduction give $k_{\mathrm{tree}}\approx4.3\times10^5\approx2^{18.7}$,
so the IT soundness error is at most $2^{-109.3}$. After compilation, the state-restoration bound is
$\kappa(t)\le(t+k_{\mathrm{tree}})\cdot2^{-128}+4.5\max(t,t^*)^2/2^{256}$
(\cref{subsec:sec-cap}). In the regime $t\ge2k_{\mathrm{tree}}$ this is at most $t/2^{127}$, so $N$
enters only through $k$. Two caveats must be written into the theorem [A2-m4]:
\begin{itemize}[nosep]
\item The theorem is stated for the sequentialised protocol with \emph{independent} challenges per
  node, which is what Fiat--Shamir per node gives. If parallel folds shared challenges, a pair state
  $S_1\wedge S_2$ would flip whenever either component flips, and the errors would add.
\item The side condition $2k\le t$ of the compiled headline applies to $k_{\mathrm{tree}}$.
\end{itemize}

\paragraph{Labels.} Until \lean{tree\_secure} is proved, an L1 record certifies the one-step
theorems only, and it says so. Even then, at L1 the extractor's running time is unconstrained. A
brute-force extractor satisfies the definition (\cref{subsec:sec-extractor}), so the tree theorem is
an IT \emph{soundness} statement. Calling a T2 record \enquote{knowledge-sound} or
\enquote{recursion-ready} requires the L2 efficient-extraction track. The reason is that the
recursive bound below consumes the extractor time $t_E$ explicitly [A2-M2].

\begin{figure}[t]
\centering
\begin{tikzpicture}[
  every node/.style={font=\small},
  leaf/.style={draw,rounded corners=2pt,fill=black!3,inner sep=2.5pt,align=center,font=\scriptsize},
  acc/.style={draw,rectangle split,rectangle split parts=2,rectangle split horizontal,
              inner sep=2pt,font=\scriptsize,rectangle split part fill={blue!8,orange!12}},
  fold/.style={draw,ellipse,inner sep=1pt,fill=blue!4,font=\scriptsize\itshape},
  cost/.style={font=\scriptsize,align=left,anchor=west,text width=4.6cm},
  arr/.style={-{Latex[length=1.5mm]}},
  qry/.style={densely dashed,-{Latex[length=1.2mm]},black!50}]
\foreach \j/\x in {1/0,2/2.5,3/5.0,4/7.5} {
  \node[leaf] (r\j) at (\x,0) {$(\vk_i,x_{\j})$\\[-1pt]$\idx(i)$ via \texttt{QUERY}};
  \node[acc] (a\j) at (\x,1.45) {$\mathrm{pl}_{\j}$\nodepart{two}$\mathcal O_{\j}$};
  \draw[arr] (r\j) -- node[right,font=\scriptsize\itshape,text=blue!50!black]{embed} (a\j);
}
\node[fold] (f12) at (1.25,2.65) {fold};
\node[fold] (f34) at (6.25,2.65) {fold};
\foreach \a/\f in {a1/f12,a2/f12,a3/f34,a4/f34} { \draw[arr] (\a) -- (\f); }
\node[acc] (b1) at (1.25,3.75) {$\mathrm{pl}_{12}$\nodepart{two}$\mathcal O_{12}$};
\node[acc] (b2) at (6.25,3.75) {$\mathrm{pl}_{34}$\nodepart{two}$\mathcal O_{34}$};
\draw[arr] (f12) -- (b1); \draw[arr] (f34) -- (b2);
\node[fold] (f) at (3.75,4.85) {fold};
\draw[arr] (b1) -- (f); \draw[arr] (b2) -- (f);
\node[acc] (c) at (3.75,5.9) {$\mathrm{pl}_{1..4}$\nodepart{two}$\mathcal O_{1..4}$};
\draw[arr] (f) -- (c);
\node[draw,rounded corners=2pt,fill=green!8,font=\scriptsize] (d) at (3.75,7.0) {decide: accept / reject (native)};
\draw[arr] (c) -- (d);
% cost column
\node[cost] at (9.0,0.75) {\textbf{embed}, $N\times$: $S_{\mathrm{leaf}}$; reads $\vk_i$, $x$
  ($n_x\le2^{10}$), queries $\idx(i)$; reduces $x$ to a claim};
\node[cost] at (9.0,3.2) {\textbf{fold}, $(N-1)\times$: $S_{\mathrm{fold}}$, \emph{the recursion
  step}; reads two plain parts ($\le2^{12}$ $K$-elements each) from the tape and queries
  $\mathcal O$; output = same shape};
\node[cost] at (9.0,6.45) {\textbf{decide}, $1\times$: $S_{\mathrm{dec}}/64$ in
  $\mathrm{Score}_{T2}$; bounded only by the fuel};
% legend
\node[acc,anchor=west] (lg) at (-0.6,-1.0) {pl\nodepart{two}$\mathcal O$};
\node[font=\scriptsize,anchor=west,align=left] at (0.35,-1.0) {accumulator: plain part on the tape $|$ $\le16$ oracles via \texttt{QUERY}};
\end{tikzpicture}
\caption{Where recursion pays. A 2-to-1 tree for $N=4$. Each accumulator has the fixed class shape:
a plain part on the tape and at most 16 committed oracles. Each arrow is a reduction with a one-step
Lean theorem (L1). The whole tree is covered by the organiser's \lean{tree\_secure} (L2, IT model).
The deployed system, where each fold verifier is compiled with a hash and proved inside the next
node's proof, rests on the named assumption \lean{RecursionHeuristic}
(\cref{subsec:rec-cannot}).}
\label{fig:rec-tree}
\end{figure}

%-----------------------------------------------------------------------------------------------
\subsection{What cannot be a Lean theorem}\label{subsec:rec-cannot}

The deployed system compiles each node's fold verifier with a hash $H$ (BCS and Fiat--Shamir) and
proves \enquote{the fold verifier accepted} inside the next node's proof. The proved relation
therefore queries $H$. Three results, summarised in \cref{subsec:bg-recursion}, bound what can be
said about that system.
\begin{enumerate}[nosep]
\item Succinct arguments for relations that themselves query the random oracle (relativized
  arguments) do not exist in the ROM, even interactive ones with computational soundness
  \cite{BCG24-RelROM}. A ROM proof of recursive soundness is therefore ruled out, not merely missing.
\item After the oracle is instantiated by a concrete hash, standard-model PCD theorems assume an
  extractor for the instantiated scheme, and they cover only constant depth: an extractor of size
  $|\tilde\Prv|^c$ per layer costs $|\tilde\Prv|^{c^d}$ at depth $d$ \cite[Rmk.~5.3]{BCMS20}. A
  2-to-1 tree has depth $\log_2N=\omega(1)$.
\item The only depth-free concrete bound \cite{CGSY24} needs a straight-line extractor for a
  relativized argument. For hash-based systems that is available only in a non-succinct idealisation,
  so it is heuristic evidence.
\end{enumerate}
leanerVM records the same gap as its open problem T6 (\enquote{recursive CTE}). It keeps the
end-to-end theorem as a target schema and forbids a global bridging axiom.\footnote{\repo{leanerVM@8aaf8a0}, \path{README.md:35-38},
\path{docs/architecture.md:297-350}, \path{AGENTS.md:81-83}.} The competition follows that policy.

\begin{assumption}{RecursionHeuristic}{Recursive deployment (named, never a Lean hypothesis of a scored theorem)}
Replacing the random oracle by the deployed hash inside the recursion circuit preserves straight-line
extraction of the accumulation scheme, with extractor time $t_E$. Under this assumption the advertised
security of a tree of $N$ nodes is $\kappa_{\mathrm{ARG}}(\lambda,q,s+O(N\cdot t_E))$, in the form
of \cite{CGSY24}. In the trusted repository this is prose plus an \emph{uninhabited} structure in
leanerVM's target-schema style:
\begin{leancode}
structure RecursionHeuristic (A : Accumulation k) (H : HashFn) : Prop where
  straightline_extractable : ...   -- prose in AUDIT.md; no instance
-- target schema; never proved, never a hypothesis of a scored theorem:
-- theorem recursive_sound (hA : A.Secure) (h : RecursionHeuristic A H)
--     (c : CompilerCorrect …) : …
\end{leancode}
\end{assumption}

Two more deployment facts are stated next to the assumption and are not part of any theorem.
\begin{itemize}[nosep]
\item \emph{Statement binding.} The R1CS is generated from programs, which is the precondition of
  \cite{Fen26-GenR1CS}; the harness's transcript prefix answers it by absorbing
  $\mathrm{MTroot}(\idx(i))$, $\vk_i$ and $x$ before the first challenge
  (\cref{subsec:sec-fs}). The recursion circuit also contains the Fiat--Shamir hash, the
  precondition of \cite{KRS25}; that attack is not answered by binding and is part of
  \lean{RecursionHeuristic}. For a fold node this means the two input plain parts and the roots of
  their accumulator oracles.
\item \emph{Extractor time.} The bound above is meaningful only if $t_E$ is polynomial. At L1 it is
  unconstrained, so the L2 efficient-extraction track is a precondition of the advertised bound.
\end{itemize}

%-----------------------------------------------------------------------------------------------
\subsection{Bounded-depth schemes and declared depth}\label{subsec:rec-depth}

Accumulation from linear codes without homomorphism \cite{BMNW24} checks linearity by spot-checks.
The new accumulator is then only guaranteed to be $\delta$-close to the claimed combination, and
every extraction step loses $\delta$. Soundness holds for a decider family
$\mathsf{D}_0\subseteq\dots\subseteq\mathsf{D}_d$, where $\mathsf{D}_s$ accepts words that are
$s\delta$-close to the code. Once $d\delta$ exceeds the unique-decoding radius there is a concrete
attack, not a gap in the proof. Arc and WARP remove the depth bound for RS and for general linear
codes respectively \cite{BMNW25-Arc,BCFW25-WARP}. The Lean shape of a bounded-depth scheme is
level-indexed: a relation $\Racc^{(s)}$ for each level $s\le d_{\max}$, a fold from
$\mathsf{pair}(\Racc^{(s)})$ to $\Racc^{(s+1)}$, and a decider per level.

\begin{decision}{REC-3}{Depth-uniform in the MVP; depth always declared and enforced}
The MVP certificate is depth-uniform (\cref{def:rec-acc}). The level-indexed generalisation is a
documented L2 extension. Every T2 record declares $d_{\max}$, and the harness and deployment
documentation enforce $d\le d_{\max}$. This applies even to depth-uniform schemes, because
\lean{tree\_secure} and the advertised bound are stated for $d\le d_{\max}$. \emph{Rationale:}
bounded-depth attacks are real, and arity bounds do not bound depth. leanVM's
\code{MAX\_RECURSIONS}$=16$ bounds the number of children per node, and leanerVM found no decreasing
depth measure in leanVM's aggregation.\footnote{\repo{leanerVM@8aaf8a0},
\path{docs/leanvm-target.md:90-94}; \repo{leanVM@a386121f},
\path{crates/rec\_aggregation/src/aggregation.rs:78}.} [D1 §4.2, §4.4; R7 §4.2]
\end{decision}

%-----------------------------------------------------------------------------------------------
\subsection{Full recursion (T3) is future work}\label{subsec:rec-t3}

The brief's \enquote{ambitious variant} arithmetises the verifier and measures a 2-to-1 step: the
fold verifier itself becomes an R1CS instance, and the step proves $\Vrf(\pi_1)\wedge\Vrf(\pi_2)$.
We schedule this as track T3, for three reasons.
\begin{enumerate}[nosep]
\item \emph{A trusted arithmetiser.} T3 needs an RV64IM+Zbc$\to$R1CS compiler whose correctness is
  part of the trusted base. That is a circuit-to-ISA correctness project of SP1/OpenVM scale. None
  exists for the contract's machine model.
\item \emph{Wall time.} A 2-to-1 wall-clock time depends on the machine and the prover
  implementation, and it is not reproducible. It can be a Yukon-style throughput board outside the
  Lean contract.
\item \emph{The fixed point does not close.} At $r$ R1CS constraints per RISC-V cycle, a verifier of
  $C$ cycles has an $r\cdot C$-constraint circuit. CPU-style arithmetisations need $r\approx30$--$100$
  (a judgement). \Cref{tab:rec-selfrec} applies this to the binary-field trajectory of
  \cref{sec:benchmark}. A self-recursive fixed point inside the $2^{24}$ class needs
  $C\le2^{24}/r$, i.e.\ $1.7$--$5.6\times10^5$ cycles. The best modelled design (B3) is 9--29 times
  above that.
\end{enumerate}

\begin{table}[t]
\centering\small
\begin{tabular}{@{}lrrr@{}}
\toprule
Verifier (cycles at the $2^{24}$ class) & $C$ & $\log_2(rC)$, $r=30$ & $r=100$ \\
\midrule
B1: Spartan + binary WHIR-JB, $\rho=1/2$ & $6.2\times10^6$ & 27.5 & 29.2 \\
B3: B1 at $\rho=1/16$ & $4.8\times10^6$ & 27.1 & 28.8 \\
B5: B1, unique-decoding regime & $1.1\times10^7$ & 28.3 & 30.0 \\
leanVM 2$\to$1 step, naive port to RV64IM+Zbc & $\approx3\times10^7$ & 29.8 & 31.5 \\
\midrule
fixed point inside $2^{24}$ requires & $\le2^{24}/r$ & 24 & 24 \\
\bottomrule
\end{tabular}
\caption{Self-recursion arithmetic over $K/E$ [D3 §6.2, reading 6; D2 §7.4]. Cycle counts are model
estimates ($\pm2\times$). leanVM itself runs a 2$\to$1 step in 0.56--0.69\,M steps of its
field-native ISA \cite{leanVM}; a leanISA step is worth 1--139 weighted RV64IM+Zbc cost units depending on the
opcode (an $E$ multiplication is 103 instructions, 139 units under v0.2's display weights with \texttt{clmul} at 4), so the two numbers are not comparable.}
\label{tab:rec-selfrec}
\end{table}

Closing the gap needs field and hash support in the ISA, which is the field-syscall board T1-F with an
\texttt{EMUL} syscall priced from a binary-field zkVM, or classes up to $2^{28}$
(\cref{sec:cost}). The contract's class ladder is nested and can be extended without changing theorem
shapes. In the meantime, T2's fold-verifier cycle count on the contract machine is the competition's
proxy for recursion cost: it is what a fixed zkVM would arithmetise. The L4 statement exists only as
the target schema \lean{recursive\_sound} above.

\begin{takeaway}[Summary]
The brief's batcher becomes an accumulation scheme: embed, fold and decide, each with the base
track's certificate, an accumulator shape fixed by type, and output agreement on the machine. Lean
certifies one step (L1) and, once the frame rule exists, the tree (L2), both in the IT model. Security
of the compiled recursive system is the named \lean{RecursionHeuristic}, with declared depth. The
fold verifier's cycles are the recursion cost that T2 scores, against the \lean{fullRecAcc} baseline
that equals T1.
\end{takeaway}
